% ===================================================================
% SKILL: SORU FRAME + DURAK FRAME + TAM COZUM FRAME
% ===================================================================
% AMAC:
% Soruyu kutusuz ana govdede vermek, gerekirse sagda gorsel kullanmak ve
% cozumu iki asamada gostermek:
% 1. dusuk kontrastli cozum icin durak
% 2. tam kirmizi cozum
%
% ON KOSUL:
% \solutionghost ve \solutiontext yardimcilari 10_icerik_boxlari.tex
% yaklasimina gore tanimlanmis olmalidir.
% ===================================================================

\begin{frame}[t]{Soru 2: Uretim Hatti}
\small
\begin{columns}[T,onlytextwidth]
    \begin{column}{0.54\textwidth}
        Bir elektronik montaj hattinda iki test asamasi vardir: termal test ve
        titresim testi.

        \vspace{0.2cm}
        \textbf{Verilenler}
        \begin{itemize}
            \item Termal testi gecme olasiligi: \(0.80\)
            \item Termal testi gecen bir modulun titresimi de gecme olasiligi: \(0.75\)
        \end{itemize}

        \textbf{Istenen}
        \begin{itemize}
            \item[A.] Rastgele secilen bir modulun iki testi de gecme olasiligini bulunuz.
            \item[B.] Sonucu kisa bir karar cumlesiyle yorumlayiniz.
        \end{itemize}
    \end{column}
    \begin{column}{0.42\textwidth}
        \centering
        \begin{tikzpicture}[scale=0.85, every node/.style={font=\scriptsize}]
            \node[draw, rounded corners, fill=blue!8, minimum width=3.2cm, minimum height=0.85cm, align=center] (t) at (2.0,2.6) {Termal testi\\gecenler};
            \node[draw, rounded corners, fill=green!10, minimum width=3.0cm, minimum height=0.85cm, align=center] (v) at (2.0,1.0) {Titresimi de\\gecen alt kume};
            \draw[-{Stealth[length=2.2mm]}, thick] (t) -- (v);
            \node[align=center] at (2.0,1.8) {kosullu okuma\\ustten alta};
        \end{tikzpicture}
    \end{column}
\end{columns}
\end{frame}

\begin{frame}[t]{Soru 2: Cozum Icin Durak}
\small
\solutionghost{
\textbf{Adim 1.} Ortak gecis olasiligi carpim mantigi ile okunur.

\[
    P(T \cap V) = P(T)\,P(V\mid T)
\]

\textbf{Adim 2.} Sayilari yerlestir:

\[
    P(T \cap V) = 0.80 \times 0.75
\]

\textbf{Adim 3.} Sonucu once kendin tamamla ve dur.
}
\end{frame}

\begin{frame}[t]{Soru 2: Cozum}
\small
\solutiontext{
\textbf{Adim 1.} Ortak gecis olasiligi carpim mantigi ile okunur:

\[
    P(T \cap V) = P(T)\,P(V\mid T)
\]

\textbf{Adim 2.} Sayilari yerlestir:

\[
    P(T \cap V) = 0.80 \times 0.75 = 0.60
\]

\textbf{Yorum.} Rastgele secilen bir modulun iki testi de gecme olasiligi
\%60'tir; bu nedenle titresim asamasi sistemde ciddi bir ikinci filtre gorevi gorur.
}
\end{frame}
